% RESULTADO_RECUPERADO: K31-10/P10 del inventario REV05.
% Fuente de las incidencias:
% 16_CIERRE_GLOBAL_HMT_MD_2026-07-22/02_LECTURA_DODECAFASICA/continuacion_fases/
% verificar_obstruccion_y_dato_minimo_w72.py, lineas 45--60 y 141--233.
% RESULTADO_OBSTRUCCION_Y_DATO_MINIMO_W72.json, phase_charge_reader.
% Fuente de las bandas: N69_signatures_t0_t99.csv, tiempos impresos abajo.
% Consumidor: PAQUETE_CONTINUIDAD_K_UNIDAD_20260918_BASE_COMPARTIDA/python/
% selector_orbital_original.py, verify_historical_inputs, lineas 159--193.
% No se transfiere la restriccion del lector comparativo al generador completo.
\subsection{Incidencias temporales del repertorio terminal y recuperación por posición}
\label{sec:repertorio-incidencias-completas}

La publicación del repertorio terminal utiliza los mismos bloques regionales,
los mismos acarreos y el mismo calendario que producen el descriptor inicial
de veinticuatro trits. Conviene explicitar por separado dos operaciones de
esta composición. La primera recupera ocho palabras a partir de incidencias
temporales registradas. La segunda presenta esas ocho palabras al selector
dodecafásico como un conjunto no ordenado. El selector vuelve a establecer
el orden mediante sus condiciones de incidencia, de cilindro y de cierre.
La distinción permite conservar los testigos de procedencia sin convertir
el orden de una segmentación histórica en una entrada del selector.

Sea $B_t=(b_{t,0},b_{t,1},b_{t,2})\in(\mathbb F_3^6)^3$ la frontera
regional ya generada en el tiempo $t$, con $0\leq t<100$. En la carta
utilizada por el registro histórico, la matriz de incidencia es
\[
 A_{\mathrm{hist}}=
 \begin{pmatrix}
 0&1&1&1&1&1\\
 1&0&1&1&2&2\\
 1&1&0&2&1&2\\
 1&1&2&0&2&1\\
 1&2&1&2&0&1\\
 1&2&2&1&1&0
 \end{pmatrix}.
\]
Las palabras se consideran vectores fila. Sus dos cargas y la fase
cronológica son
\[
 r_{t,i}=\sum_{j=1}^{6}b_{t,i,j},\qquad
 d_{t,i}=\sum_{j=1}^{6}(b_{t,i}A_{\mathrm{hist}})_j,
 \qquad \vartheta_t=t-1\pmod 3.
\]
Todas las operaciones de estas tres expresiones tienen lugar en
$\mathbb F_3$. El subíndice histórico identifica una carta concreta de
la incidencia de Witt; mantiene explícita la transformación de coordenadas
cuando se utiliza otra presentación de esa misma incidencia.

Para $c\in\{r_t,d_t\}$, $\delta\in\mathbb F_3$ y
$\epsilon\in\{1,-1\}$, definimos
\[
 \eta_i(c,\delta,\epsilon)=
 \begin{cases}
  \epsilon,&c_i=\vartheta_t+\delta,\\
  -\epsilon,&c_i\ne\vartheta_t+\delta,
 \end{cases}
 \qquad
 L_{t,c,\delta,\epsilon}
   =\sum_{i=0}^{2}\eta_i(c,\delta,\epsilon)b_{t,i}.
\]
La orientación $-1$ se escribe como $2$ al imprimir los coeficientes
ternarios. Junto a estos doce lectores comparativos por frontera, el
registro dispone del lector afín
\[
 L^{\mathrm{af}}_t
    =\sum_{i=0}^{2}(d_{t,i}+\vartheta_t)b_{t,i}.
\]
La adición de la fase en el lector afín y la comparación de cargas en
el lector anterior son operaciones distintas, ambas sobre datos ya
producidos por la frontera TPK.

\begin{table}[htbp]
\centering\small
\begin{tabular}{rccc cc}
\toprule
$t$&$b_{t,0}$&$b_{t,1}$&$b_{t,2}$&$r_t$&$d_t$\\
\midrule
8 &200121&212212&110110&011&202\\
11&022212&110001&100021&001&012\\
30&210112&110202&211110&100&012\\
34&121022&000011&112110&220&021\\
37&120121&011202&112221&100&201\\
49&110212&200212&100122&110&201\\
58&111121&022121&122201&122&220\\
67&010012&001220&011120&122&122\\
75&220000&020111&100002&120&021\\
81&122201&020202&201122&202&001\\
90&022112&022110&121010&202&200\\
\bottomrule
\end{tabular}
\caption{Fronteras temporales que realizan las incidencias impresas del
repertorio terminal. Cada entrada de seis cifras conserva el orden de
las seis coordenadas de la frontera regional.}
\label{tab:terminal-bandas-testigo}
\end{table}

La tabla siguiente despliega todas las incidencias comparativas del
registro sobre los doce bloques de la segmentación considerada. Las
letras $r$ y $d$ designan respectivamente la carga visible y la carga
dual. La columna $j$ conserva el índice de posición del testigo histórico;
la operación de entrega al selector olvida este índice después de formar
el conjunto de palabras.

\begin{table}[htbp]
\centering\small
\begin{tabular}{rrcrrc c}
\toprule
$j$&$t$&carga&$\delta$&$\epsilon$&$(\eta_0\eta_1\eta_2)$&$L_{t,c,\delta,\epsilon}$\\
\midrule
4&11&$d$&0& 1&212&021101\\
5&67&$d$&1&$-1$&211&002001\\
5&67&$d$&2& 1&211&002001\\
5&67&$r$&1&$-1$&211&002001\\
5&67&$r$&2& 1&211&002001\\
6&81&$d$&0& 1&222&020111\\
6&81&$r$&2& 1&222&020111\\
7&34&$r$&1&$-1$&111&200110\\
7&75&$d$&1&$-1$&211&200110\\
7&75&$r$&2&$-1$&211&200110\\
8&90&$r$&0& 1&121&121012\\
8&90&$r$&1&$-1$&121&121012\\
9&49&$d$&0&$-1$&121&010122\\
11&11&$d$&2&$-1$&211&221110\\
11&37&$d$&0&$-1$&121&221110\\
12&8&$d$&0&$-1$&111&222110\\
12&8&$r$&1&$-1$&111&222110\\
12&58&$d$&1&$-1$&111&222110\\
12&58&$r$&0&$-1$&111&222110\\
\bottomrule
\end{tabular}
\caption{Incidencias fase--carga completas sobre la segmentación terminal
en las cien fronteras del registro. Las coincidencias de palabras entre
filas conservan sus tiempos, orientaciones y tipos de carga.}
\label{tab:terminal-incidencias-completas}
\end{table}

El décimo bloque dispone de un testigo afín. En $t=30$ se tiene
$d_{30}=(0,1,2)$ y $\vartheta_{30}=2$, de donde
\[
 (d_{30,0}+\vartheta_{30},d_{30,1}+\vartheta_{30},
 d_{30,2}+\vartheta_{30})=(2,0,1)
\]
y, usando las bandas de la tabla~\ref{tab:terminal-bandas-testigo},
\[
 L^{\mathrm{af}}_{30}
 =2(210112)+0(110202)+(211110)
 =(001001)\quad\text{en }\mathbb F_3^6.
\]
La igualdad es coordenada a coordenada: antes de reducir módulo tres,
el vector de la suma es $(6,3,1,3,3,4)$.

\begin{proposition}[Recuperación íntegra del repertorio por incidencias]
\label{prop:repertorio-terminal-ocho-incidencias}
Considérese el registro de incidencias comparativas de la
tabla~\ref{tab:terminal-incidencias-completas}, aumentado con el testigo
afín $(j,t,L)=(10,30,001001)$. Para cada posición $5\leq j\leq12$, el
conjunto de palabras de salida de sus incidencias es unitario. La
extracción de esas ocho palabras y el olvido posterior de su orden
producen exactamente
\[
 \mathcal S_{\mathrm{term}}=
 \{002001,020111,200110,121012,010122,001001,221110,222110\}.
\]
Su codificación como enteros de base tres es
\[
 \{28,55,98,175,437,498,687,714\}.
\]
\end{proposition}
\begin{proof}
La multiplicidad de incidencias por posición $j=5,6,7,8,9,10,11,12$
es, respectivamente,
\[
 (4,2,3,2,1,1,2,4).
\]
En cada posición, todas las filas imprimen la misma palabra. Cada fila
se verifica sustituyendo sus tres bandas y los tres coeficientes en la
fórmula de $L$; el décimo caso se ha calculado explícitamente mediante
$L^{\mathrm{af}}_{30}$. Por ejemplo, en $t=67$ y con coeficientes
$(2,1,1)$ se obtiene
\[
 2(010012)+(001220)+(011120)=(002001)\pmod3.
\]
Los ocho resultados son distintos como palabras de longitud seis. Su
codificación $\sum_{k=1}^{6}w_k3^{6-k}$ da, en el orden de las
posiciones, $(55,175,498,437,98,28,714,687)$. Al olvidar el orden se
obtiene el conjunto de enteros enunciado. Así, la extracción preserva
cada palabra completa, incluidas sus posiciones nulas, y entrega ocho
elementos distintos al selector.
\end{proof}

La prueba anterior es una recuperación exacta desde incidencias
identificadas. Su dato de entrada comprende el registro de incidencias
con sus posiciones y testigos. La selección orbital posterior recibe
únicamente $\mathcal S_{\mathrm{term}}$ y el descriptor regional inicial;
recorre las $8!$ ordenaciones y determina su compatibilidad con las
fronteras funcionales y las incidencias orientadas de la rueda
dodecafásica. Ambas operaciones quedan así contiguas, con dominios
explícitos: procedencia temporal de las ocho palabras y determinación
constructiva de su orden.

El cuadro comparativo contiene las diecinueve incidencias que su fórmula
produce sobre la segmentación de doce bloques en los cien tiempos
declarados. Esa exhaustividad finita se comprueba recorriendo
$100\cdot2\cdot3\cdot2=1200$ evaluaciones. El lector afín amplía la
cobertura de las posiciones terminales de siete a ocho. Este recuento
caracteriza el repertorio local de lectores que acaba de definirse;
las demás operaciones del TPK conservan los dominios y funciones
constructivas establecidos en sus respectivas derivaciones.
